\section{Four Callers, One Query}
\label{sec:ch15-four-callers-one-query}
\index{authorization!what a refusal looks like|(}%

With a scope on \code{Customer.email} and a router that will accept a token,
one query is enough to see the whole behaviour. It is the catalogue page from
chapter~\ref{ch:hotchocolate-quickly}, extended one field:

\begin{minted}[fontsize=\footnotesize,breaklines]{graphql}
{ products { title reviews { edges { node { author { displayName email } } } } } }
\end{minted}

Four callers ask it. Nothing changes between them except the
\code{Authorization} header.

\begin{minted}[fontsize=\footnotesize,breaklines]{text}
no token             200   one error per protected field instance; author null
token, read:pii      200   nothing refused, email present
token, no scope      200   one error per protected field instance; author null
token, expired       401   one error, no data at all
\end{minted}

\subsection*{The two that look alike and are not}

\index{UNAUTHORIZED\_FIELD\_OR\_TYPE@\code{UNAUTHORIZED\_FIELD\_OR\_TYPE}}%
The anonymous caller and the scopeless caller get the same shape and the same
error code, \code{UNAUTHORIZED\_FIELD\_OR\_TYPE}, and differ in one word:

\begin{minted}[fontsize=\scriptsize,breaklines]{text}
Unauthorized to load field 'Query.products.reviews.edges.node.author.email', Reason: not authenticated.

Unauthorized to load field 'Query.products.reviews.edges.node.author.email', Reason: missing required scopes.
\end{minted}

That is a good error message. It names the full path rather than the field, so
a client developer can find it in their own document, and the reason
distinguishes \enquote{log in} from \enquote{you are logged in and this is not
yours}, which are different things to show a user.

\index{authorization!a null that is not a refusal}%
The second row of that table says \enquote{nothing refused} rather than
\enquote{no errors}, and the difference cost me a gate run. Asked with the
right scope on a freshly seeded database, the query answers with no errors at
all. Asked at the end of a full verification pass it answers with one, because
that pass submits a review against a customer key that was never seeded, and a
query that walks every author in the catalogue eventually walks that one. The
author comes back null and the response carries an error, and neither has
anything to do with a token.

So the gate asserts that no error carries
\code{UNAUTHORIZED\_FIELD\_OR\_TYPE}, rather than that there are no errors.
Chapter~\ref{ch:strangling-the-monolith} met this same trap from the other
side and stopped a resolver count short of \code{Product.reviews} for exactly
this reason: a number, or an assertion, that holds only on a database nothing
has written to is true once. It is also the practical shape of the distinction
this section is about, because from the client's side an authorisation null and
a resolution null look identical until you read the error beside them.

\subsection*{A denied field takes its parent with it}

\index{nullability!and a denied field}%
The line in that table I did not expect is \enquote{author null}. I removed one
field and lost the object it was on.

The reason is nullability and it has been sitting in this book since
chapter~\ref{ch:entity-resolution-done-right}. \code{email} is \code{String!}.
The router cannot put null in a non-null field, so the null propagates to the
nearest nullable ancestor. \code{Review.author} is nullable, and that is where
it stops.

It is nullable because of an argument
chapter~\ref{ch:strangling-the-monolith} made, and that argument had nothing to
do with authorisation: a reference to an entity another subgraph owns must be
nullable, because the router may fail to find the thing. Had that argument gone
the other way, one denied \code{email} would have taken the entire response to
\code{data: null}, and a catalogue page would have gone blank because one
reviewer's address was private.

So boundary nullability turns out to decide the blast radius of an
authorisation rule as well as the blast radius of a failed fetch. Nobody aimed
it there. If you are putting \code{@requiresScopes} on a field today, look up
the chain from it and find where the first nullable ancestor is, because that
is what your users lose.

\subsection*{Two doors, two spellings}

\index{Cosmo Router!401 versus a field error}%
The expired token is refused somewhere else entirely. It never reaches
planning: the router's access controller rejects it, the response is HTTP 401,
and the body is:

\begin{minted}[fontsize=\footnotesize,breaklines]{text}
{"errors":[{"message":"unauthorized"}]}
\end{minted}

with no \code{data} key at all. A bad token is a fact about the request. A
missing scope is a fact about one field of one response. They deserve different
answers and they get them, but a client that branches on
\enquote{unauthorised} as a single condition will handle one of these two
correctly and the other one wrongly, and which one depends on whether the
author was thinking about logging in or about permissions that day.

\subsection*{The two switches Mosaic leaves alone}

\index{authorization!require\_authentication}%
The \code{authorization} block has three keys and all three default to false.
Mosaic ships the first two false and I turned each on once to see what the
graph becomes.

\code{require\_authentication: true} makes the router answer 401 to any request
with no token, before planning, whatever it asked for. Asking only for product
titles, anonymously:

\begin{minted}[fontsize=\footnotesize,breaklines]{text}
HTTP 401
{"errors":[{"message":"unauthorized"}]}
\end{minted}

That is the switch that turns a public graph into an internal one, and it is a
single line. It is also why Mosaic leaves it off: this is a storefront, and a
storefront that demands a token before it will list a product does not sell
anything.

\index{authorization!reject\_operation\_if\_unauthorized}%
\code{reject\_operation\_if\_unauthorized: true} changes the field decision
into an operation decision. With a scopeless token the whole query now fails,
and something new appears in the response:

\begin{minted}[fontsize=\scriptsize,breaklines]{text}
{"errors":[{"message":"Unauthorized"}],"data":null,"extensions":{"authorization":
{"missingScopes":[{"coordinate":{"typeName":"\0\0\0\0\0\0\0\0","fieldName":"email"},
"required":[["read:pii"]]}],"actualScopes":[]}}}
\end{minted}

HTTP 200, not 401, which is the third status this one condition has now
produced. Note also \code{Unauthorized} with a capital here against the
lowercase \code{unauthorized} of the 401 path: two code paths, two spellings,
one word.

The \code{extensions.authorization} block is genuinely useful and I would want
it in development. It says which coordinate was refused, what that coordinate
required, and what the caller actually had, which is the three facts you need
to work out whose fault it is.

Choosing between the two is a product decision rather than a security one.
Neither leaks the value. One gives a client a partial answer it has to be
written to handle; the other gives it a clean failure and nothing to render.
Mosaic keeps the partial answer, because a catalogue page missing one address
is still a catalogue page.

\subsection*{The eight zero bytes}

\index{Cosmo Router!zeroed typeName in missingScopes}%
The \code{typeName} in that listing should say \code{Customer}. On the wire it
is eight NUL bytes. The listing above writes them as \code{\textbackslash 0}
eight times because this book will not print a control character. The
substitution is mine; the bytes are real.

It reproduced three times out of three, identically. \code{Customer} is eight
characters long, so the length is right and the content is not, which is what a
string built at the correct size and never filled looks like. The
\code{fieldName} beside it is correct.

The router does not build that value. Its authorizer is handed a graph
coordinate by the execution engine, which is \code{graphql-go-tools} v2.14.2,
and records what it was given. I did not read the engine, so I am not going to tell you where the
zeroing happens. What I can tell you is that it reaches a client, and one
consequence that is not obvious from looking at it: \code{addMissingScopes}
deduplicates by comparing type name and field name, so with every type name
zeroed, two protected fields called the same thing on two different types would
be reported once.

This is the second defect this book has found in a released version of one of
these tools, after the composer crash chapter~\ref{ch:hard-modeling-problems}
left in the open items, and it goes to the same place: worth reporting
upstream, and nobody has. It is not in a gate, and I want to be straight about
why. The setting that produces it is off in the shipped configuration, so
asserting the bytes would mean a verification run that turns a switch on to
catch a defect in somebody else's release. It is recorded in the open items
instead, with the exact edit that reproduces it.

\medskip

Every refusal so far has come from the router. The other field, the one the
service was supposed to protect, has not been asked for yet.

\index{authorization!what a refusal looks like|)}%
